\chapter{Análisis de seguridad en redes IoT} % Análisis de seguridad en redes IoT
\hyphenation{In-te-re-se}
Llegados a este punto, y utilizando como base toda la información recopilada sobre las redes IoT y los protocolos estudiados, se llevará a cabo el análisis de seguridad sobre las redes IoT, con el objetivo de conocer sus vulnerabilidades y debilidades en el diseño, funcionamiento e implementación. A partir de este análisis, se establecerán una serie de conclusiones sobre los riesgos de seguridad existentes en las redes IoT.




\section{Auditoría de seguridad en una red industrial} % Auditoría de seguridad en una red industrial
Se procede a realizar una auditoría de seguridad Wireless en la red Zigbee del laboratorio creado (Apéndice \ref{laboratorio}). Esta auditoría será de caja blanca ya que se conoce la configuración de la red y los dispositivos que la componen, simulando que ha sido llevada a cabo por un auditor de seguridad contratado para evaluar la red.
\subsection{Planificación} % Planificación
Para llevar a cabo la auditoría se seguirá el esquema y los controles de la metodología OWISAM, descritos al principio del proyecto. De esta forma, se llevarán a cabo de forma ordenada las siguientes acciones:
\begin{itemize}
\item Recopilar toda la información sobre la red y los dispositivos que la forman.
\item Analizar las comunicaciones entre dispositivos.
\item Analizar los mecanismos de cifrado de las comunicaciones.
\item Verificar la configuración de la red.
\item Llevar a cabo pruebas sobre la infraestructura y su disponibilidad.
\item Llevar a cabo pruebas sobre los dispositivos.
\end{itemize}
\clearpage




\subsection{Recopilación de información e Identificación de dispositivos} % Recopilación de información e Identificación de dispositivos
La primera fase a llevar a cabo es la recopilación de información sobre el sistema a auditar, mediante técnicas pasivas (recopilación de información sobre el objetivo sin interacción directa: monitorización del tráfico de red) y activas (recopilación de información sobre el objetivo con interacción directa: escaneo de red). La recopilación de toda esta información simplifica las siguientes fases.
\\\\El objetivo de esta fase es conocer en profundidad la configuración de la red, el flujo de información entre dispositivos y todas las tecnologías implicadas en la comunicación. Para llevar a cabo esta fase, se han seguido tres estrategias:
\begin{itemize}
\item Capturar tráfico de la red Zigbee mediante la placa de sniffing Api-Mote y el framework Killerbee. Especificamente las herramienta zbstumbler para encontrar el canal de Zigbee en funcionamiento y zbdump para capturar el tráfico IEEE 802.15.4
\item Capturar tráfico de la red Wifi con la herramienta Wireshark.
\item Realizar varios escaneos en la red con la herramienta Nmap.
\end{itemize}
\vspace{0.6cm}
En la siguiente tabla se muestra la información recopilada más relevante. Dicha información sobre la configuración de la red y los dispositivos, junto con el análisis de las capturas de tráfico realizadas, permite a un atacante conocer la composición de la red, ya que se exponen los flujos de datos/información que se dan entre los dispositivos y los diferentes protocolos que engloban el conjunto de la red formada por el usuario.

\clearpage
\begin{table}[H]
\begin{adjustwidth}{-7em}{-4.5em}
\begin{center}
\begin{tabular}{|l|l|}
\hline
\rowcolor[HTML]{91B9F5} 
\multicolumn{1}{|c|}{\cellcolor[HTML]{91B9F5}\textbf{Información}}                                                                    & \multicolumn{1}{c|}{\cellcolor[HTML]{91B9F5}\textbf{Método/Indicio}}                                               \\ \hline
Red Zigbee basada en IEEE St 802.15.4-2003                                                                                            & Campo ''Frame version'' en captura de tráfico                                                                 \\ \hline
\begin{tabular}[c]{@{}l@{}}Modo de direccionamiento de red: direcciones\\ cortas (16 bits)\end{tabular}                               & \begin{tabular}[c]{@{}l@{}}Campos ''Source/Destination Addressing mode''\\ en captura de tráfico\end{tabular} \\ \hline
\begin{tabular}[c]{@{}l@{}}Seguridad a nivel de red y aplicación\\ habilitada (los mensajes entre dispositivos \\viajan cifrados)\end{tabular}                                            & \begin{tabular}[c]{@{}l@{}}Campos ''Security y Security Enabled'' en\\ captura de tráfico\end{tabular}        \\ \hline
\begin{tabular}[c]{@{}l@{}}Identificador de red de área personal \\ (PAN ID): 0x3fa5\end{tabular}                                     & Campo ''Destination PAN'' en captura de tráfico                                                               \\ \hline
\begin{tabular}[c]{@{}l@{}}Identificador extendido de red de área\\ personal (EPID): d0:d9:d8:dd:6b:81:c5:e7\end{tabular}             & Campo ''Extended PAN ID'' en captura de tráfico                                                               \\ \hline
Dirección corta de la puerta de enlace: 0x0000                                                                                        & Campo ''Source'' en captura de tráfico                                                                        \\ \hline
\begin{tabular}[c]{@{}l@{}}Dirección extendida de la puerta de enlace: \\ ec:1b:bd:ff:fe:0d:25:f0\end{tabular}                       & Campo ''Extended Source'' en captura de tráfico                                                               \\ \hline
\begin{tabular}[c]{@{}l@{}}Direcciones cortas de sensores Zigbee: 0xf1e7,\\ 0xced2\end{tabular}                                       & Campo ''Source'' en captura de tráfico                                                                        \\ \hline
\begin{tabular}[c]{@{}l@{}}Direcciones extendidas de sensores Zigbee:\\ 14:b4:57:ff:fe:3c:9d:76, 14:b4:57:ff:fe:16:e8:22\end{tabular} & Campo ''Extended Source'' en captura de tráfico                                                               \\ \hline
\begin{tabular}[c]{@{}l@{}}Aplicación de control de la red: TuyaSmart\end{tabular}                                                                                         & Captura de tráfico Wifi \\ \hline
\begin{tabular}[c]{@{}l@{}}Ip y Mac de la puerta de enlace Zigbee: \\ 192.168.1.41 / 68:57:2D:31:D7:1D \\(Hangzhou Aixiangji Technology)\end{tabular}                                                                                         & Escaneo de hosts activos en la red                                                                                                                                                                                                                                              \\ \hline
\begin{tabular}[c]{@{}l@{}}Ip y Mac de la cámara ip: \\ 192.168.1.42 / 68:57:2D:4B:E8:D2 \\(Hangzhou Aixiangji Technology)\end{tabular}                                                                                                & Escaneo de hosts activos en la red                                                                                                                                                               \\ \hline
\begin{tabular}[c]{@{}l@{}}Puerto y servicio de comunicación entre \\ Wifi y Zigbee: 6668/TCP - Protocolo IRC\end{tabular}            & Escaneo de puertos abiertos y servicios en el host \\ \hline
\begin{tabular}[c]{@{}l@{}}Ip y Mac del Smartphone que controla \\la red: 192.168.1.40 / EC:10:7B:89:08:53 \\(Samsung Electronics)\end{tabular} & Escaneo de hosts activos en la red                                                                                \\ \hline
\end{tabular}
\end{center}
\end{adjustwidth}
\caption{Recopilación de información en la red Zigbee}
\end{table}


\clearpage
\subsection{Ataques a la red Zigbee} % Ataques a la red Zigbee
En esta fase se ejecutarán algunos de los ataques analizados (Apéndice \ref{ataquesredesieee}) contra la red Zigbee. Los ataques efectuados se dividen en las siguientes categorías:
\subsubsection{Ataques de repetición} % Ataques de repetición
Los ataques de repetición/playback son un tipo de ataques en los cuales un atacante intercepta información de la red, para retransmitirla posteriormente, con el fin de suplantar la identidad de uno de los lados de la comunicación o impedir la misma.
\\\\
Zigbee cuenta con una contramedida ante este tipo de ataques, el número de secuencia de las tramas. Este contador se va incrementando con cada paquete válido que se transmite a la red, siendo el dispositivo receptor el encargado de verificar y desechar las tramas que contengan un contador incorrecto.
\\\\
Sin embargo, en el caso de estudio de este trabajo donde se analizan los dispositivos del fabricante Ksix, se ha descubierto una vulnerabilidad desconocida \cite{exploitdb} que permite evadir el mecanismo de protección y cuya explotación posibilita que un atacante falsifique todo tipo de mensajes entre dispositivos de forma remota.
\\\\
Esta vulnerabilidad es producida por una debilidad en el software de los dispositivos Zigbee de tipo CWE-294 (''authentication bypass by capture-replay''), por la cual al capturar un mensaje de la red, modificando el campo ''número de secuencia'' por un valor elevado como 250, se evade el mecanismo de protección al transmitir de nuevo el mensaje en un momento específico de la comunicación. Esta vulnerabilidad fue reportada a Ksix y la PoC fue aceptada y publicada en Exploit-DB.
\\\\
Durante esta fase de auditoría, se ha explotado la vulnerabilidad en la red Zigbee desarrollando un script en bash, el cual modifica el valor del campo ''número de secuencia'' de un paquete previamente capturado que contiene el mensaje ''abrir puerta'', retransmitiéndolo a la red utilizando la herramienta zbreplay del framework killerbee, y provocando que el mensaje falsificado sea analizado y aparezca una alerta en la aplicación de control de la red que tiene el usuario.
\clearpage
\subsubsection{Ataques de denegación de servicio} % Ataques de denegación de servicio
La ejecución de este tipo de ataques en la red Zigbee impiden el normal funcionamiento de la misma, provocando la pérdida de mensajes y conectividad.
\\\\
Para realizar estos ataques se han utilizado las herramientas zbrealign, forzando a un dispositivo final a asociarse con otra red, y zbassocflood, enviando masivamente solicitudes de unión al coordinador (inundación de paquetes).
\\\\
Por medio de la inundación de paquetes, se ha conseguido dejar inoperativa la red Zigbee por un intervalo de tiempo suficiente (aproximadamente 2 minutos) para que algunas de las alertas generadas por los dispositivos en ese espacio de tiempo se hayan perdido y otras hayan llegado desfasadas. Pasado ese tiempo el coordinador rechaza los paquetes manipulados y el ataque deja de tener efecto.
\\\\A contraposición con el tipo de ataque de denegación de servicio anterior, el envío de un paquete de asociación manipulado a un dispositivo final para se conecte a otra red, que puede o no existir, consigue una denegación total de las funciones del dispositivo. El dispositivo queda totalmente inoperativo sin que el usuario se dé cuenta ya que todavía aparece como conectado en la aplicación de control de la red.
\vspace{1cm}
\subsubsection{Ataques criptográficos} % Ataques criptográficos
Se han realizado ataques a la confidencialidad de las comunicaciones entre los dispositivos Zigbee sin éxito, comprobando que se protege adecuadamente la confidencialidad en la red. El objetivo principal de estos ataques era la adquisición de las claves de cifrado (claves de enlace y red) pudiendo posteriormente descifrar los mensajes enviados entre dispositivos y analizar los datos en texto plano.
\\\\Mediante las capturas de tráfico Zigbee y desarrollando un pequeño script en bash que utiliza la herramienta zbdsniff, se ha realizado un ataque de diccionario y fuerza bruta verificando que no se utilizan las contraseñas por defecto que define Zigbee y que se utiliza el estándar AES con claves de 128 bits, longitud suficiente para considerar ''seguras'' las claves utilizadas.

\clearpage


\section{Medidas de protección a ataques contra redes IoT} % Medidas de protección a ataques contra redes IoT
Realizado el estudio de los múltiples vectores de ataque que existen en una red de IoT Zigbee y llevada a cabo la auditoría de seguridad, se establecen en este punto una serie de medidas que sirvan de referencia para realizar de forma segura la implementación y configuración de una red IoT:
\begin{itemize}
\item Configurar un sistema de detección de intrusiones inalámbrico (WIDS). Los WIDS pueden capturar conductas maliciosas que logren vulnerar la primera capa de defensa de una red, es decir, ataques contra la criptografía y la autenticación.
\item Permitir la conexión a la red únicamente a dispositivos autorizados, descartando cualquier tráfico procedente de un dispositivo no conocido. Para llevar a cabo esto, se pueden definir listas blancas o negras con la información de los dispositivos que tendrán permitido o denegado el acceso.
\item Utilizar varios enrutadores, configurando una topología de red en malla, consiguiendo que la red no se sature fácilmente.
\item Limitar el número de peticiones de acceso a la red, estableciendo un límite máximo de intentos consecutivos erróneos.
\item Implementar mecanismos para el correcto restablecimiento de la red en caso de saturación y fallo debido a la ejecución de un ataque DoS.
\item Implementar el cifrado de las comunicaciones en la red mediante algoritmos y claves robustas. Se recomienda el uso de AES utilizando claves de un mínimo de 128 bits.
\item Implementar mecanismos para validar el número contador de secuencia que eviten la recepción y procesamiento de paquetes repetidos.
\item Diseñar un esquema de validación y comprobación de mensajes duplicados en el nivel de aplicación, asegurando así la integridad y el no repudio.
\item Implementar marcas de tiempo en los mensajes enviados, evitando el procesamiento de estos en un periodo de tiempo posterior al definido.
\item Controlar el acceso físico a los dispositivos y proteger el firmware del dispositivo cifrando el mismo.
\item Configurar el arranque seguro del dispositivo, comprobando la integridad del firmware y evitando arrancar si existe una modificación en los datos.
\item Implementar dispositivos IoT con la funcionalidad de resetearse a valores de fábrica y borrar parámetros sensibles, ante la detección de una apertura del dispositivo o acceso indebido. Esta funcionalidad se denomina reducción a cero.
\end{itemize}



\clearpage



